% ch-terms.tex -- theory: the individual physics terms of the flux
\chapter{The physics terms}
\label{ch:terms}

\Cref{eq:flux} collects six flux contributions, and \cref{eq:height} adds
evaporation and a source. Each of them has its own sub-block under
\key{Physics} and its own \key{Enable} key, so they are switched on and off one
at a time. This chapter takes them in turn: what the term is, what it does to
the film, what behavior you should expect to see when it dominates, and which
keys control it. The keys are quoted by name only --- their types, defaults and
ranges are tabulated once, in \cref{ch:ref-physics}, and are not repeated here.

\section{Capillary (Laplace) pressure}
\label{sec:terms:laplace}

\begin{equation}\label{eq:q-cap}
  \Qvec_{\mathrm{cap}}
    = \frac{\hat{\sigma}(C)\,H^{3}}{3\hat{\mu}(C)}\,\nabla\nabla^{2}(H+S).
\end{equation}

Surface tension acting to flatten the free surface. It is responsible for
almost everything the word ``leveling'' covers: a rough coating smoothing
itself out, a drop relaxing toward a spherical cap, liquid collecting in the
dips of a textured substrate. In practice nearly every run keeps it on.

Two details are worth holding on to. First, \textbf{this term is why the height
equation is fourth order}: $\nabla\nabla^{2}$ is three derivatives acting on
$H$ inside a divergence, so \cref{eq:height} carries a fourth spatial
derivative, and that in turn is why each boundary needs two conditions
(\cref{sec:conc:boundaries}) and why the linear systems are what they are
(\cref{ch:linearsolver}). Second, the curvature is measured on the free surface
$H+S$, not on the film thickness $H$ --- which is how the substrate shape gets
into the problem at all.

Its mobility contains $\hat{\sigma}(C)$ as well as $\hat{\mu}(C)$, so as a film
dries and its surface tension changes, the rate at which it levels changes too.

The only key is \key{Physics.LaplacePressure.Enable}
(\cref{sec:ref-physics:laplace}).

\section{Disjoining pressure}
\label{sec:terms:disjoining}

\begin{equation}\label{eq:q-dis}
  \Qvec_{\mathrm{dis}} = \Dpn\,\frac{H^{3}}{3\hat{\mu}}\,\nabla\Pihat(H),
  \qquad
  \Pihat = (\Hstar/H)^{n} - (\Hstar/H)^{m},
  \qquad n > m > 1.
\end{equation}

Once a film becomes very thin, the substrate and the free surface start to
interact with one another. Expressing that interaction as a pressure gives the
disjoining pressure, and it is this term that gives the model partial wetting:
a finite contact angle and a contact line that can move. The first term of
$\Pihat$ is repulsive and tries to keep the film thick; the second is
attractive. They cancel at $H = \Hstar$, where $\Pihat = 0$, and that thickness
is the precursor film left behind outside a drop.

The characteristic behavior is a drop of finite extent sitting on a thin
precursor, with a transition region between the two whose width is set by
$\Hstar$ and the contact angle. How the coefficient $\Dpn$ is calibrated
against the angle you ask for, and how wide that transition region actually is,
is the subject of \cref{ch:wetting}.

Four keys live in the \key{Physics.DisjoiningPressure} block: \key{Enable},
\key{ContactAngle}, \key{ExponentN} and \key{ExponentM}. The precursor
thickness comes from \key{Initial.SurfaceShape.PrecursorRatio}
(\cref{sec:ref-physics:disjoining}). Enabling the term with
\key{PrecursorRatio} left at \val{0} leaves $\Hstar$ undefined and terminates
the run.

\begin{warn}[This is the stiffest term in the model]
  $\Pihat \sim H^{-n}$ is a very strong non-linearity: with the conventional
  exponents $(n,m) = (3,2)$ a halving of the thickness multiplies the repulsive
  branch by 8. They are the defaults of \key{ExponentN} and \key{ExponentM}
  (\cref{sec:ref-physics:disjoining}).
  Nothing else in the equation responds that steeply to a small change in $H$.
  Expect this term to dominate the Newton iteration, the conditioning of the
  linear system and the admissible time step whenever it is on.
\end{warn}

\section{Vertical gravity and the substrate shape}
\label{sec:terms:vgravity}

\begin{equation}\label{eq:q-vg}
  \Qvec_{\mathrm{vg}} = -\frac{\Bo}{3}\,\frac{H^{3}}{\hat{\mu}}\,\nabla(H+S).
\end{equation}

Gravity normal to the substrate. Hydrostatic pressure drives liquid from the
high parts of the free surface to the low ones, so on a horizontal substrate it
flattens bumps and works alongside capillary pressure; the two together are
what level a film. The difference between them is the order of the derivative,
and hence the length scale: gravity levels long-wavelength unevenness that
capillary pressure barely feels.

The keys are \key{Physics.VerticalGravity.Enable} and
\key{Physics.\bk VerticalGravity.\bk Acceleration}, the latter in
\si{\metre\per\second\squared} (\cref{sec:ref-physics:vgravity}). If you are
using a substrate tilted by an angle $\alpha$, the value to give here is
$g\cos\alpha$, and the in-plane component goes to
\cref{sec:terms:hgravity}.

\subsection*{The substrate shape}

The substrate profile $S(x,y)$ is built in the \key{Substrate} block. Any
number of \key{RectBump}, \key{EllipseBump}, \key{PeriodicWave},
\key{FlatPlane}, \key{MexicanHat} and \key{SmoothConvex} blocks may be listed;
they are applied in ascending \key{Order} and combined into one shape
(\cref{sec:ref-domain:substrate}). The substrate never moves and never changes
with time.

$S$ enters the equations only through the combination $H+S$, and only in
\cref{eq:q-cap,eq:q-vg}. \textbf{If neither capillary pressure nor vertical
gravity is enabled, the substrate shape does nothing whatsoever} --- it is
parsed, built, written to the output and then never read by any term. This is
a common way to spend an afternoon wondering why a carefully built topography
has no effect.

\section{In-plane gravity}
\label{sec:terms:hgravity}

\begin{equation}\label{eq:q-hg}
  \Qvec_{\mathrm{hg}} = \frac{\Pvec}{3}\,\frac{H^{3}}{\hat{\mu}}.
\end{equation}

Flow down an inclined substrate. What distinguishes it from vertical gravity is
that the driving field is a constant vector, independent of position: there is
no gradient in it at all. The flux is then simply proportional to $H^{3}$, so
\textbf{thicker parts of the film run faster}, the thick regions overtake the
thin ones, and a front forms at the leading edge. That front steepening is the
signature behavior of this term, and it is the origin of the fingering
instabilities seen in real coatings running downhill.

The keys are \key{Physics.HorizontalGravity.Enable} and
\key{Physics.\bk HorizontalGravity.\bk Acceleration}, a pair $(g_x, g_y)$ in
\si{\metre\per\second\squared} (\cref{sec:ref-physics:hgravity}). For a
substrate tilted by $\alpha$, give the in-plane component $g\sin\alpha$ ---
\num{4.903} at $30^\circ$. The sign sets the downhill direction.

\section{Centrifugal force}
\label{sec:terms:centrifugal}

\begin{equation}\label{eq:q-cf}
  \Qvec_{\mathrm{cf}} = \frac{\Spn}{3}\,\frac{H^{3}}{\hat{\mu}}\,r\,\er.
\end{equation}

A rotating substrate throws the film outwards. This is the main driving force
of spin coating. The driving field is proportional to the distance $r$ from the
rotation axis, so unlike in-plane gravity it is not uniform: the effect is
weak at the center and strongest at the rim, which is what produces the
characteristic thinning-from-the-middle history of a spin-coated film.

The keys are \key{Physics.Centrifugal.Enable},
\key{Physics.Centrifugal.RotationRate} in \si{\radian\per\second}, and
\key{Physics.Centrifugal.RotationAxis}, the coordinates of the rotation center
in meters (\cref{sec:ref-physics:centrifugal}).

\begin{note}[What a 1-D spin-coating run does with the $y$ direction]
  In a \val{"xh"} run there is no $y$ direction to be thrown into, but the
  liquid would be thrown there all the same. The solver accounts for it: the
  material that would leave through the unsolved $y$ direction is treated as
  loss, $\partial_y[\tfrac13\mathrm{Sp}\,H^3(y-y_0)/\hat\mu] = \tfrac13\mathrm{Sp}\,H^3/\hat\mu$
  --- the line $y = y_0$ of a film that is uniform in $y$ there. A 1-D
  spin-coating run is therefore not a slice through a 2-D one, and its total
  volume falls even between walls. Both models do this; \val{"CZ"} also
  carries the sideways divergence into its vertical velocity. What leaves this
  way is counted as outflow in \file{History.dat} (\code{OutflowedVol},
  \code{OutflowedResin}), so \code{VolErr} still balances
  (\cref{sec:output:history}).
\end{note}

A rotating film is not periodic: the arm $r$ grows across the domain and jumps
by its length across a periodic seam, where liquid would be created or lost.
The solver refuses \val{"Periodic"} edges with the centrifugal term
(\cref{sec:ref-bcout:boundary}); a strip of a film that goes on beyond the domain
takes an \val{"Outlet"} at each end. The one exception is a run with a
manufactured solution (\key{Numerics.ManufacturedSolution}), whose source is
formed for the arm of each cell: the periodic domains of \file{tools/mms.py}
carry every term at once.

\section{Solutal capillarity (Marangoni)}
\label{sec:terms:marangoni}

\begin{equation}\label{eq:q-sc}
  \Qvec_{\mathrm{sc}} = \frac{\As\,\Ma}{2}\,\frac{H^{2}}{\hat{\mu}}\,\nabla C.
\end{equation}

A surface stress arising because surface tension depends on concentration.
Where drying is faster the concentration rises, surface tension changes with
it, and the resulting gradient drags liquid along the free surface. Drying
mottle and flow streaks in a coating commonly start here.

Note the power of $H$. \textbf{This term scales with $H^{2}$, not $H^{3}$,
because it is a shear applied at the surface rather than a body force acting
through the whole thickness.} Integrating a constant body force across a
parabolic profile gives $H^{3}/3$; integrating the response to a surface
traction gives $H^{2}/2$. One consequence is that Marangoni flow becomes
relatively more important as a film thins, which is precisely when
concentration gradients are strongest.

The keys are \key{Physics.SolutoCapillarity.Enable} and
\key{Fluid.SurfaceTension.Coef}, which is $\mathrm{d}\sigma/\mathrm{d}c$ in
\si{\newton\per\metre} per unit concentration
(\cref{sec:ref-physics:soluto}). The sign of \key{Coef} decides which way the
liquid is pulled. Setting \key{Enable = false} forces $\Ma = 0$, which removes
this term and, in the same stroke, removes the concentration dependence of
surface tension everywhere else (\cref{sec:conc:properties}).

\section{Navier slip at the substrate}
\label{sec:terms:slip}

Slip is not a term of its own but a change of every mobility. With a slip
length $\lambda$ the velocity at the substrate is $\lambda\,\partial u/\partial z$
instead of zero, the parabolic profile gains a uniform slip velocity, and the
flux integrated across the film becomes
%
\begin{equation}\label{eq:q-slip}
  \Qvec = -\frac{1}{\mu}\Big(\frac{h^{3}}{3} + \lambda h^{2}\Big)\nabla p
          + \frac{1}{\mu}\Big(\frac{h^{2}}{2} + \lambda h\Big)\nabla\sigma .
\end{equation}
%
In the dimensionless equations $H^{3}/3$ of \cref{eq:q-cap,eq:q-dis,eq:q-vg,eq:q-hg,eq:q-cf}
becomes $H^{3}/3 + \Lambda H^{2}$ and $H^{2}$ of \cref{eq:q-sc}
becomes $H^{2} + 2\Lambda H$, with $\Lambda = \lambda/h_0$. The extra part grows
relative to the no-slip one as the film thins, $3\Lambda/H$ against 1: slip
matters most at a contact line, where it relieves the stress singularity as a
precursor film does by other means. It leaves the driving potential, and
therefore every equilibrium, unchanged. The keys are in
\cref{sec:ref-physics:slip}; the well-mixed model only.

\section{Evaporation}
\label{sec:terms:evaporation}

Evaporation is not a flux contribution. It removes solvent only, so it enters
the film-height equation and not the balance of the solute:

\begin{yzntable}{Where the evaporation rate enters\label{tab:evap-terms}}{@{}l L{0.5\textwidth}@{}}
\textsf{\bfseries Equation} & \textsf{\bfseries Evaporation term} \\ \midrule
Height, \cref{eq:height} & $-\Ev\,\hat{E}(H,C)$ \\
Solute $HC$, \cref{eq:conc} & none \\
Concentration $C$, the form of \cref{eq:conc-adv} & $+\Ev\,\hat{E}(H,C)\,C/H$ \\
\end{yzntable}

The film thins and the solute left behind becomes more concentrated. Written for
$C$, the enrichment is divided by $H$, so thin regions concentrate fastest ---
the same observation made about \cref{eq:conc-adv}.

The rate factorizes into three independent pieces, plus a correction:

\begin{equation}\label{eq:evap-factor}
  \hat{E} = \mathrm{Coef}\cdot f_{\mathrm{conc}}(C)\cdot
            f_{\mathrm{space}}(X)\cdot f_h(H).
\end{equation}

\begin{description}
  \item[Concentration, $f_{\mathrm{conc}}$] chosen by \key{ConcFunc}.
    \val{"Raoult"} is Raoult's law: the rate is proportional to the mole
    fraction of solvent, computed from the molar-mass ratio $\Mr$ of
    \key{Fluid.MolarMass}. \val{"PowerN"} is
    $(1 - C - \mathrm{ConcR})^{\mathrm{PowerN}}$, floored at zero so it never
    goes negative.
  \item[Position, $f_{\mathrm{space}}$] chosen by \key{SpatialFunc}.
    \val{"uniform"} is the same everywhere in the plane. \val{"negPow"} is a
    power of the distance to the rim of a circular pool of radius
    $R = \mathrm{NegPowRmax}$ centered on the origin of the grid coordinates,
    \begin{equation}\label{eq:evap-negpow}
      f_{\mathrm{space}} = \min\!\Bigl(k\,\bigl(1 - r/R'\bigr)^{-n},\ 10\Bigr),
      \qquad k = \tfrac{1}{2}(1-n)(2-n),\quad R' = 1.001\,R,
    \end{equation}
    with $r = \sqrt{x^2+y^2}$ and $n = \mathrm{NegPowExponent} < 1$; $k$
    makes the mean over the disc $r < R'$ equal to 1, which needs $n < 1$.
    For $0 < n < 1$ the rate diverges toward the rim, capped at 10; $n = 0$
    is uniform, and for $n < 0$ the rate falls to zero at the rim. Over the last 5 cells before $R$ the curve is blended
    into a quarter cosine that takes over its value and slope and levels off
    at $R$, and the rate is constant beyond $R$.
  \item[Precursor correction, $f_h$] a correction that suppresses evaporation
    from the precursor film. It is applied only when \key{PrecursorRatio} is
    greater than zero, and it ramps up linearly: \textbf{zero at $H = \Hstar$,
    and 1 at $H = 20\Hstar$.} It is zero below $\Hstar$ as well, so a film
    thinner than the precursor neither evaporates nor takes up solvent. Its purpose is to stop the precursor drying out
    and disappearing, so it is always active in a calculation that contains a
    contact line.
\end{description}

No step removes more solvent than the film holds: $E\Delta t \le h(1-c)$, at
either end of the step and in both equations, beside the cap at
$h/(2\Delta t)$. A film therefore dries toward $c = 1$ without passing it ---
with a rate that does not depend on $c$, a third of the remaining solvent is
left after each step --- and no solute is ever removed. Raoult's law takes $c$
within $[0, 1]$.

\key{Physics.Evaporation.Coef} has the dimensions of a velocity,
\si{\metre\per\second}, and is the source of both $\Ev$ and $\Pen$ --- changing
it moves the balance against diffusion as well as the drying rate. The
condition under which a run is declared dry and stopped is set by
\key{EndConcType} (\val{"mean"}, \val{"max"} or \val{"min"}) and
\key{EndConcValue}. All of these keys are tabulated in
\cref{sec:ref-physics:evaporation}.

\section{The source (nozzle)}
\label{sec:terms:source}

The supply term $\hat{S}$ added to \cref{eq:height}. It is modeled as a nozzle
mounted on the end of a rotating arm dispensing liquid onto the film, which
covers both a stationary dispense and a moving one. The keys live in
\key{Physics.Source} (\cref{sec:ref-physics:source}).

\begin{description}
  \item[Footprint] \key{Radius} is the radius of the nozzle mouth in meters. A
    cell receives supply in proportion to the fraction of it covered by that
    circle --- by that interval, in 1-D. \key{SubRegions} sets how many
    sub-divisions per cell are used to measure the fraction, so a larger value
    makes the supply vary more smoothly as the nozzle moves. Each
    sub-division counts when its center lies inside the footprint.
  \item[Rate] \key{InletRate} is the steady-state supply, given with the
    dimensions of a rate of thickness increase, \si{\metre\per\second}.
    \key{TimeConstant} is the rise-time constant $\tau$ in seconds, with the
    rate growing as $1 - e^{-t/\tau}$.
  \item[Arm] \key{ArmCenter} is the pivot position in meters and
    \key{ArmLength} the arm length, also in meters; \key{RotationRate} is its
    angular velocity in \si{\radian\per\second} and \key{StartAngle} the
    initial angle in radians.
    \key{ArmLength = 0} gives a nozzle that does not move. \key{RampTime} and
    \key{PlateauTime} are the ramp time applied to the rotation rate and the time
    it holds its plateau: the rate is $\omega f(t)$, with $f$ rising linearly
    from 0 to 1 over $T_s$, holding 1 for $T_m$ and falling back to 0 over
    another $T_s$, and the arm's angle is
    $\theta = \theta_0 + \omega\int_0^t f\,dt'$, so the arm stops at
    $\theta_0 + \omega(T_s + T_m)$ and stays there.
\end{description}

\paragraph{The concentration of the supplied liquid.}
\key{ConcType} says what the arriving liquid brings to the solute balance
\cref{eq:conc-dim}, as $c_S\,S$ (in the well-mixed model; \val{"CZ"} ignores
the nozzle altogether, \cref{sec:ref-physics:disjoining}):
\begin{description}
  \item[\val{"local"} (the default)] The liquid has the concentration already
    present where it lands, $c_S = c$. It thickens the film without changing
    $c$, and \key{Conc} is not used.
  \item[\val{"supplied"}] The liquid has the concentration \key{Conc}, a mass
    fraction in $[0:1]$, $c_S = \key{Conc}$. It is mixed through the
    thickness of the column it lands on, as the well-mixed model assumes of
    all its liquid, so the solute $hc$ of the column grows by $S\,c_S$ while
    its film grows by $S$. On a film that does not flow,
    \begin{equation}\label{eq:source-mixing}
      c(t) = \frac{h_0 c_0 + c_S\int_0^t S\,dt'}{h_0 + \int_0^t S\,dt'},
    \end{equation}
    so a thin film under the nozzle takes on $c_S$ quickly and a thick one
    slowly; in general the concentration is pulled toward $c_S$ at the rate
    $S/h$ (\cref{eq:conc-adv}).
\end{description}
\file{History.dat} counts the solute that came with the liquid in
\code{SuppliedResin} (\cref{sec:output:history}), at the concentration the
mode gives it.

